(1) 定态方程 $-\hbar^2\psi''/(2m)+V(x)\psi=E\psi$。对于 $-V_0<E<0$，阱内波数 $k=\sqrt{2m(E+V_0)}/\hbar$ 为实数，阱外衰减常数 $\kappa=\sqrt{-2mE}/\hbar>0$。阱内一般解为 $A\cos kx+B\sin kx$，阱外可归一化解必须沿 $|x|\to\infty$ 指数衰减，不能选发散的指数项。
(2) 对称势允许本征态取偶或奇宇称。偶态阱内取 $A\cos kx$，外侧取 $A\cos(ka)e^{-\kappa(|x|-a)}$，对 $x=a$ 匹配导数得 $k\tan(ka)=\kappa$。奇态阱内为 $A\sin kx$，外侧带 $\operatorname{sgn}x$，匹配给 $-k\cot(ka)=\kappa$。并有 $k^2+\kappa^2=2mV_0/\hbar^2$，三个条件共同决定离散能级。
(3) 定义无量纲阱强 $\rho=a\sqrt{2mV_0}/\hbar$，在第一象限画 $z=ka$ 与 $\sqrt{\rho^2-z^2}$ 的交点。偶态方程 $z\tan z=\sqrt{\rho^2-z^2}$ 在任意 $\rho>0$ 的 $0<z<\min(\rho,\pi/2)$ 内至少有一根。奇态首次出现于束缚能恰接近零、$\kappa\to0$，此时 $-\cot z=0$ 的最小正根为 $z=\pi/2$，故需 $\rho>\pi/2$ 才有真正负能奇束缚态。
(4) 偶态的概率归一化为 $1=A^2[\int_{-a}^{a}\cos^2(kx)\,dx+2\cos^2(ka)\int_a^\infty e^{-2\kappa(x-a)}dx]$，所以 $A^{-2}=a+\sin(2ka)/(2k)+\cos^2(ka)/\kappa$。当 $E\uparrow0$ 时 $\kappa\to0$，阱外衰减长度 $\ell_{\rm out}=1/\kappa$ 发散，粒子概率大量延伸到阱外；这也是浅束缚态容易受长程扰动影响的原因。
